\section{Hipótesis de partida}
\label{sec:hypotheses}

Este trabajo parte de una hipótesis principal sobre la arquitectura completa y
la descompone en afirmaciones contrastables para las capas estratégica, mecánica,
operacional y de red. Cada una se evalúa dentro de los supuestos del subproblema
correspondiente, mediante análisis formal, experimentos pareados o ablaciones. Por
tanto, una propiedad verificada en una capa no se extiende automáticamente al
sistema integrado.

\subsection{Hipótesis principal}

\noindent\textbf{HP.}
Cuando la flota disponga de recursos suficientes, la configuración Cargo sea
factible en el plano y el horizonte sea compatible con la duración de la tarea,
una arquitectura híbrida que combine decisiones estratégicas locales,
comunicación vecinal, cierres enteros y guardias físicas explícitas formará
coaliciones factibles y completará la misión de transporte dentro de una región
de operación medible. La comparación con referencias globales cuantificará la
pérdida de calidad y el coste de la información. El apoyo a HP exige identificar
una región reproducible de éxito y confirmar, mediante ablaciones, la contribución
de sus bloques críticos. El enunciado no presupone estabilidad global del sistema
híbrido, una ejecución enteramente distribuida ni validez en hardware.

\subsection{Hipótesis específicas}

\noindent\textbf{H1a. Realizabilidad estratégica.}
En instancias realizables de SP0 y SP1, una penalización mayor que el coste
unilateral máximo hará coincidir los equilibrios de Nash puros con asignaciones
de cobertura o cuota exacta. Un equilibrio deficitario que satisfaga los
supuestos formales constituirá un contraejemplo.

\noindent\textbf{H1b. Cierre entero.}
En condiciones de escasez, el cierre por ranking y cuórum (QR) aplicado a la
relajación de Smith reducirá la fracción de robots que permanece en coaliciones
parciales. La evidencia no apoyará esta afirmación si el efecto pareado frente a
Smith continuo es nulo o adverso.

\noindent\textbf{H1c. Incentivo de cuórum.}
Con el cierre QR fijado, el beneficio de cuórum aumentará el valor de las cargas
completadas frente al beneficio lineal. Solo recibirá apoyo si la comparación
pareada favorece al incentivo de cuórum.

\noindent\textbf{H2. Factibilidad mecánica.}
La guardia vectorial de \emph{wrench} planar reducirá las falsas aceptaciones
respecto al criterio escalar, sin una pérdida sistemática de cobertura útil en los
mundos evaluados. La afirmación quedará sin apoyo si los falsos positivos no
disminuyen o si la abstención del filtro impide su uso operativo.

\noindent\textbf{H3. Acoplamiento espacial.}
Bajo el régimen de red degradada fijado, incorporar el coste espacial a la
decisión estratégica reducirá el tiempo de terminación frente a la variante
desacoplada. Para sustentar esta ventaja, el intervalo de confianza del efecto
temporal deberá excluir cero en la dirección favorable. Una tasa de éxito
coincidente, por sí sola, no demostrará equivalencia.

\noindent\textbf{H4. Reparación.}
Tras un fallo simple, el juego restaurará el certificado aditivo cuando la reserva
completa cubra el déficit. En la misión integrada, activar la reparación aumentará
el éxito frente a desactivarla. La hipótesis perderá apoyo si alguna instancia
recuperable no restaura el certificado o si la diferencia de éxito no es
favorable; el agotamiento del horizonte se registrará como fallo de recuperación.

\noindent\textbf{H5a. Coste de escala.}
En el catálogo binario de SP8, el coste de una activación local crecerá con el
grado vecinal, mientras que el del oráculo exhaustivo implementado crecerá con el
número de perfiles. La separación deberá observarse dentro del dominio en el que
el oráculo certifica su solución.

\noindent\textbf{H5b. Calidad bajo escala.}
La tasa global de perfiles libres de conflicto del método con retransmisión
periódica no colapsará al aumentar la escala en el dominio ensayado. La hipótesis
se refutará si la tasa alcanza cero en algún tamaño o si su deterioro contradice
la conservación postulada.

Los contrastes estocásticos emplearán mundos y semillas pareados. El análisis
conservará las colisiones, los fallos de convergencia y los agotamientos del
horizonte. Las pruebas de SP5 y SP7 se interpretarán como evidencia local sobre
seguridad y tráfico, no como una garantía global del sistema integrado.
